\begin{afterword}
\summaryhead

The post argues that the Second Law of Thermodynamics is at bottom a law about knowledge, and that the same law governs minds.

It begins with the first law: no assembly of parts can create energy if no single part can. It then explains the Second Law through Liouville's theorem, which says that the physics of a closed system never shrinks a region of possible states. A toy model with two small systems shows one system's uncertainty shrinking while the other's grows. Entropy appears to increase because we stop tracking where the states went.

The author concludes that entropy describes our beliefs about a system. If you knew the exact state of a glass of hot water, it would be cold in the thermodynamic sense, since Maxwell's demon could then extract work from it. So gaining knowledge requires observation and thermodynamic work. The post ends by saying that a clever argument cannot produce knowledge of the unseen unless some step in it breaks the laws of physics.

\responsehead

The post uses about 2,400 words of physics to reach one sentence in bold: to form accurate beliefs about something, you have to observe it. No one needs thermodynamics to accept that. The physics is there to give the sentence the standing of a law of nature, and then to support a debating rule: the reader may disbelieve any ``big, elaborate clever argument'' unless its author names the step that breaks physics.

The physics is weakest at the two points on which the conclusion depends. The first is the claim that knowing the state of hot water makes it ``colder''. Learning where the molecules are does not slow any of them. What changes is how much work the knower could extract, and the post redefines temperature as that quantity without saying so. The second is the claim that gaining information costs work. The post attributes the demon's failure to the cost of ``inspecting'' molecules. That was Szilard's and Brillouin's account, and Bennett replaced it in 1982: measurement can in principle be free, and the unavoidable cost comes from resetting the demon's memory. The post mentions that account in one parenthesis and dismisses it as ``a matter of words''. Its own demon table supports Bennett: the cost appears as the demon's memory left in one of four states.

The general rules are also stated more broadly than they hold. ``The Second Law of Thermodynamics is a corollary of Liouville's Theorem'' is taken back one paragraph later, where entropy increase becomes an ``appearance'' produced by describing the system coarsely. The rule that one subsystem narrows only if another widens fails in the post's own demon example, where the correlation between two systems is used up and no subsystem widens.

The central idea, that entropy measures an observer's information, is E.~T. Jaynes's, published in 1957. The post presents it with ``Yes, really'' and credits Jaynes only with a proverb, slightly reworded from his \textsc{Probability Theory} (2003) and given without a source.

In short: a physics lecture that gets its physics wrong at the two points on which the conclusion depends, used to present a truism as a law of nature and a debating rule as a theorem.
\end{afterword}
